\appendix
%%%%%%%%%%%%%%%%%%%%%%%%%%%%%%%%%%%%%%%%%%%%%%%%%%
%	Chapter 1 Appendix
%%%%%%%%%%%%%%%%%%%%%%%%%%%%%%%%%%%%%%%%%%%%%%%%%%
\chapter{Chapter 1 Appendix}
\graphicspath{{C:/Users/david/coding/02_TRAFFIC_BIAS_paper/6000_LaTeX/figures/}}
\addcontentsline{lof}{chapter}{\numberline{\thechapter}Chapter 1 Appendix}
\addcontentsline{lot}{chapter}{\numberline{\thechapter}Chapter 1 Appendix}

%%%%%%%%%%%%%%%%%%%%%%%%%%%%%%%%%%%%%%%%%%%%%%%%%%
%	Scrutiny Based Event Studies 
%%%%%%%%%%%%%%%%%%%%%%%%%%%%%%%%%%%%%%%%%%%%%%%%%%
\subsection{Scrutiny Based Event Studies}
In what should again be cautioned as suggestive evidence due to concerns with endogeneity in event scrutiny inferences, I provide the following event study figures and tables of estimates below. 

%	Full Event Study Figures: Volumes
\begin{figure}
	\begin{center}
		\includegraphics[width=5.5in]{n_stops_main_full.eps}
		\caption[Total Stop Volume Event Studies by Scrutiny]{This figure plots OLS coefficients with a 95\% confidence interval for two equations, in three columns. The left column depicts coefficient estimates for all treated departments, with each point representing a 14-day bin relative to the event date, as described in equation \ref{eq2_estudy}. The right two columns come from equation \ref{eq4_estudy_int}, with the center column depicting 14-day bin estimates for low-scrutiny events, and the right column for high-scrutiny events. Standard errors are clustered at the department-level. As depicted, these figures suggest that the decrease may be driven by low scrutiny events.}
		\label{n_stops_main_full}
	\end{center}
\end{figure}

\begin{figure}
	\begin{center}
		\includegraphics[width=5.5in]{n_bstops_main_full.eps}
		\caption[Black Stop Volume Event Studies by Scrutiny]{As described in Figure \ref{n_stops_main_full}, this figure plots black stop volumes both overall and by scrutiny level, suggesting that the decrease may be driven by low scrutiny events.}
		\label{n_bstops_main_full}
	\end{center}
\end{figure}

\begin{figure}
	\begin{center}
		\includegraphics[width=5.5in]{n_searches_main_full.eps}
		\caption[Total Search Volume Event Studies by Scrutiny]{As described in Figure \ref{n_stops_main_full}, this figure plots total search volumes both overall and by scrutiny level, suggesting that the decrease may be driven by low scrutiny events.}
		\label{n_searches_main_full}
	\end{center}
\end{figure}

\begin{figure}
	\begin{center}
		\includegraphics[width=5.5in]{n_bsearches_main_full.eps}
		\caption[Black Search Volume Event Studies by Scrutiny]{As described in Figure \ref{n_stops_main_full}, this figure plots black search volumes both overall and by scrutiny level, suggesting that the decrease may be driven by low scrutiny events.}
		\label{n_bsearches_main_full}
	\end{center}
\end{figure}

\begin{figure}
	\begin{center}
		\includegraphics[width=5.5in]{n_hits_main_full.eps}
		\caption[Total Hit Volume Event Studies by Scrutiny]{As described in Figure \ref{n_stops_main_full}, this figure plots total hit volumes both overall and by scrutiny level, suggesting that the decrease may be driven by low scrutiny events.}
		\label{n_hits_main_full}
	\end{center}
\end{figure}

\begin{figure}
	\begin{center}
		\includegraphics[width=5.5in]{n_bhits_main_full.eps}
		\caption[Black Hit Volume Event Studies by Scrutiny]{As described in Figure \ref{n_stops_main_full}, this figure plots black hit volumes both overall and by scrutiny level, suggesting that the decrease may be driven by low scrutiny events.}
		\label{n_bhits_main_full}
	\end{center}
\end{figure}

%	Full Event Study Figures: Shares
\begin{figure}
	\begin{center}
		\includegraphics[width=5.5in]{bstop_share_main_full.eps}
		\caption[Black Share of Stops by Scrutiny]{As described in Figure \ref{n_stops_main_full}, this figure plots black stop-shares both overall and by scrutiny level, with little evidence of a significant difference across race.}
		\label{bstop_share_main_full}
	\end{center}
\end{figure}

\begin{figure}
	\begin{center}
		\includegraphics[width=5.5in]{bsearch_share_main_full.eps}
		\caption[Black Share of Searches by Scrutiny]{As described in Figure \ref{n_stops_main_full}, this figure plots the black search-shares both overall and by scrutiny level, with no evidence of a racially differential change.}
		\label{bsearch_share_main_full}
	\end{center}
\end{figure}

\begin{figure}
	\begin{center}
		\includegraphics[width=5.5in]{bhit_share_main_full.eps}
		\caption[Black Share of Hits by Scrutiny]{As described in Figure \ref{n_stops_main_full}, this figure plots black hit-shares both overall and by scrutiny level, with suggestive evidence of a large decrease, but only under high scrutiny.}
		\label{bhit_share_main_full}
	\end{center}
\end{figure}

%	Full Event Study Figures: Rates 
\begin{figure}
	\begin{center}
		\includegraphics[width=5.5in]{bsearch_rate_main_full.eps}
		\caption[Black Search Rate by Scrutiny]{As described in Figure \ref{n_stops_main_full}, this figure plots black search rates both overall and by scrutiny level, with suggestive evidence of an increase under high scrutiny.}
		\label{bsearch_rate_main_full}
	\end{center}
\end{figure}

\begin{figure}
	\begin{center}
		\includegraphics[width=5.5in]{bhit_rate_main_full.eps}
		\caption[Black Hit Rate by Scrutiny]{As described in Figure \ref{n_stops_main_full}, this figure plots black hit rates both overall and by scrutiny level, with suggestive evidence of a decrease under high scrutiny.}
		\label{bhit_rate_main_full}
	\end{center}
\end{figure}

\begin{figure}
	\begin{center}
		\includegraphics[width=5.5in]{bhit_bstop_ratio_main_full.eps}
		\caption[Black Hits-to-Stops Ratio by Scrutiny]{As described in Figure \ref{n_stops_main_full}, this figure plots the black hits-to-stops ratios both overall and by scrutiny level, with no evidence of an effect.}
		\label{bhit_bstop_ratio_main_full}
	\end{center}
\end{figure}

%%%%%%%%%%%%%%%%%%%%%%%%%%%%%%%%%%%%%%%%%%%%%%%%%%
%	TWO-WAY MUNDLAK DEVICE 
%%%%%%%%%%%%%%%%%%%%%%%%%%%%%%%%%%%%%%%%%%%%%%%%%%
\clearpage
\subsection{Two-Way Mundlak Device}
In order to further assess heterogeneity in treatment effects, I utilize a variation of the Two-Way Mundlak (TWM) device, from Wooldridge (2021). In my context with high frequency data and no clearly defined cohorts, it is difficult to evaluate heterogeneity. Further, in my main specification I already have some degree of insulation against these concerns as well, as I separate observations into department-events, and I use balancing to equalize observation weights across departments. Later in the appendix, I also consider log outcomes, yielding similar results, suggesting that large departments are not driving my results too strongly, using the following equations: 

\begin{equation}
    	\begin{split}
		Y^{o}_{d,t} = \beta Post_{d,t} + \boldsymbol{X'}_{d,t} \boldsymbol{\gamma} + \boldsymbol{\bar{X}}_{ym} \boldsymbol{\theta}_1 + \boldsymbol{\bar{X}}_{dy} \boldsymbol{\zeta}_2 \\ 
            + Post_{d,t} (\boldsymbol{\bar{X}}_{ym} - \boldsymbol{\bar{X}}) + Post_{d,t} (\boldsymbol{\bar{X}}_{d} - \boldsymbol{\bar{X}}) + u_{d,t} 
    	\label{eq5_TWM_DiD}
	\end{split}
\end{equation}

\begin{equation} 
    	\begin{split}
    		Y^{o}_{d,t} = \beta_1 Post^{L}_{d,t} + \beta_2 Post^{H}_{d,t} + \boldsymbol{X'}_{d,t} \boldsymbol{\gamma} + \boldsymbol{\bar{X}}_{ym} \boldsymbol{\zeta}_1 + \boldsymbol{\bar{X}}_{d} \boldsymbol{\zeta}_2 \\ 
            + Post_{d,t} (\boldsymbol{\bar{X}}_{ym} - \boldsymbol{\bar{X}}) + Post_{d,t} (\boldsymbol{\bar{X}}_{d} - \boldsymbol{\bar{X}}) + u_{d,t} 
    	\label{eq6_TWM_DiDiD}
    	\end{split}
\end{equation}

%	Full DiD + DiDiD Tables
\begin{table}
	\begin{center}
		\input{C:/Users/david/coding/02_TRAFFIC_BIAS_paper/6000_LaTeX/tables/DiD_volumes_TWM_full.tex}
		\caption[Two-Way Mundlak Volume Outcome Estimates]{This table gives volume outcome coefficient estimates for the treatment effect following a shooting event both overall using equation \ref{eq1_DiD} in Panel A and by scrutiny level using equation \ref{eq3_DiDiD} in Panel B. Despite different estimated magnitudes, the sign of these estimates generally align with my main results.}
		\label{full_DiD_volumes_TWM}
	\end{center}
\end{table}

\begin{table}
	\begin{center}
		\input{C:/Users/david/coding/02_TRAFFIC_BIAS_paper/6000_LaTeX/tables/DiD_rates_TWM_full.tex}
		\caption[Two-Way Mundlak Rate Outcome Estimates]{This table gives rate outcome coefficient estimates for the treatment effect following a shooting event both overall using equation \ref{eq1_DiD} in Panel A and by scrutiny level using equation \ref{eq3_DiDiD} in Panel B. Again, the signs of these estimates generally align with my main results.}
		\label{full_DiD_rates_TWM}
	\end{center}
\end{table}

As mentioned in the main body of this paper, these results are quite different in magnitude, but the signs of the estimates remain fairly consistent, providing some reassurance about the validity of my main results. \\

%%%%%%%%%%%%%%%%%%%%%%%%%%%%%%%%%%%%%%%%%%%%%%%%%%
%	FULL UNBALANCED SAMPLE 
%%%%%%%%%%%%%%%%%%%%%%%%%%%%%%%%%%%%%%%%%%%%%%%%%%
\clearpage
\subsection{Full Unbalanced Sample}
In efforts to utilize the largest possible sample of events and traffic stops, I have also run my regressions using a larger, unbalanced panel. Again, my results here remain largely similar. \\

%	Full DiD + DiDiD Tables
\begin{table}
	\begin{center}
		\input{C:/Users/david/coding/02_TRAFFIC_BIAS_paper/6000_LaTeX/tables/DiD_volumes_A00_full.tex}
		\caption[Volume Outcome Estimates, Full Unbalanced]{This table gives volume outcome coefficient estimates for the treatment effect following a shooting event both overall using equation \ref{eq1_DiD} in Panel A and by scrutiny level using equation \ref{eq3_DiDiD} in Panel B.}
		\label{full_DiD_volumes_A00}
	\end{center}
\end{table}

\begin{table}
	\begin{center}
		\input{C:/Users/david/coding/02_TRAFFIC_BIAS_paper/6000_LaTeX/tables/DiD_rates_A00_full.tex}
		\caption[Rate Outcome Estimates, Full Unbalanced]{This table gives rate outcome coefficient estimates for the treatment effect following a shooting event both overall using equation \ref{eq1_DiD} in Panel A and by scrutiny level using equation \ref{eq3_DiDiD} in Panel B.}
		\label{full_DiD_rates_A00}
	\end{center}
\end{table}

%%%%%%%%%%%%%%%%%%%%%%%%%%%%%%%%%%%%%%%%%%%%%%%%%%
%	DEPARTMENT-EVENT FIXED EFFECTS 
%%%%%%%%%%%%%%%%%%%%%%%%%%%%%%%%%%%%%%%%%%%%%%%%%%
\clearpage
\subsection{Department-Event Fixed Effects}
To evaluate whether the stronger restriction of department-event-year fixed effects used in my main specification are driving results, I run my regressions again using weaker department-event fixed effects, finding similar results. \\

%	Full DiD + DiDiD Tables
\begin{table}
	\begin{center}
		\input{C:/Users/david/coding/02_TRAFFIC_BIAS_paper/6000_LaTeX/tables/DiD_volumes_A01_full.tex}
		\caption[Volume Outcome Estimates, Dept.-Event FEs]{This table gives volume outcome coefficient estimates for the treatment effect following a shooting event both overall using equation \ref{eq1_DiD} in Panel A and by scrutiny level using equation \ref{eq3_DiDiD} in Panel B.}
		\label{full_DiD_volumes_A01}
	\end{center}
\end{table}

\begin{table}
	\begin{center}
		\input{C:/Users/david/coding/02_TRAFFIC_BIAS_paper/6000_LaTeX/tables/DiD_rates_A01_full.tex}
		\caption[Rate Outcome Estimates, Dept.-Event FEs]{This table gives rate outcome coefficient estimates for the treatment effect following a shooting event both overall using equation \ref{eq1_DiD} in Panel A and by scrutiny level using equation \ref{eq3_DiDiD} in Panel B.}
		\label{full_DiD_rates_A01}
	\end{center}
\end{table}

%%%%%%%%%%%%%%%%%%%%%%%%%%%%%%%%%%%%%%%%%%%%%%%%%%
%	DAY-OF-WEEK FIXED EFFECTS 
%%%%%%%%%%%%%%%%%%%%%%%%%%%%%%%%%%%%%%%%%%%%%%%%%%
\clearpage
\subsection{Day-of-Week Fixed Effects}
Noting that some departments and/or drivers may behave differently and exhibit different patterns and treatment responses across different days of the week, I have run a robustness check with the inclusion of day-of-week fixed effects as well, with similar results. \\

%	Full DiD + DiDiD Tables
\begin{table}
	\begin{center}
		\input{C:/Users/david/coding/02_TRAFFIC_BIAS_paper/6000_LaTeX/tables/DiD_volumes_A02_full.tex}
		\caption[Volume Outcome Estiamtes, DOW FEs]{This table gives volume outcome coefficient estimates for the treatment effect following a shooting event both overall using equation \ref{eq1_DiD} in Panel A and by scrutiny level using equation \ref{eq3_DiDiD} in Panel B.}
		\label{full_DiD_volumes_A02}
	\end{center}
\end{table}

\begin{table}
	\begin{center}
		\input{C:/Users/david/coding/02_TRAFFIC_BIAS_paper/6000_LaTeX/tables/DiD_rates_A02_full.tex}
		\caption[Rate Outcome Estimates, DOW FEs]{This table gives rate outcome coefficient estimates for the treatment effect following a shooting event both overall using equation \ref{eq1_DiD} in Panel A and by scrutiny level using equation \ref{eq3_DiDiD} in Panel B.}
		\label{full_DiD_rates_A02}
	\end{center}
\end{table}

%%%%%%%%%%%%%%%%%%%%%%%%%%%%%%%%%%%%%%%%%%%%%%%%%%
%	LOG OUTCOMES 
%%%%%%%%%%%%%%%%%%%%%%%%%%%%%%%%%%%%%%%%%%%%%%%%%%
\clearpage
\subsection{Log Outcomes}
To address concerns that my results are disproportionately-driven by large departments, I check log-transformed volume outcomes, with largely unchanged results. \\

%	Full DiD + DiDiD Tables
\begin{table}
	\begin{center}
		\input{C:/Users/david/coding/02_TRAFFIC_BIAS_paper/6000_LaTeX/tables/DiD_volumes_A03_full.tex}
		\caption[Volume Outcome Estimates, Log Transformation]{This table gives volume outcome coefficient estimates for the treatment effect following a shooting event both overall using equation \ref{eq1_DiD} in Panel A and by scrutiny level using equation \ref{eq3_DiDiD} in Panel B.}
		\label{full_DiD_volumes_A03}
	\end{center}
\end{table}

%%%%%%%%%%%%%%%%%%%%%%%%%%%%%%%%%%%%%%%%%%%%%%%%%%
%	STRONGER BALANCE 
%%%%%%%%%%%%%%%%%%%%%%%%%%%%%%%%%%%%%%%%%%%%%%%%%%
\clearpage
\subsection{Stronger Balance}
Checking whether my results are driven by my somewhat-weak balance definition, I run my regressions using a much stronger definition of balance, with largely unchanged results. Specifically, my main specification defines an event as balanced if I have at least one department-date observation in of the +/- 14 periods surrounding an event, whereas here I use the stronger restriction that they must have observations on \emph{every} day within +/- 14 periods. \\

%	Full DiD + DiDiD Tables
\begin{table}
	\begin{center}
		\input{C:/Users/david/coding/02_TRAFFIC_BIAS_paper/6000_LaTeX/tables/DiD_volumes_A04_full.tex}
		\caption[Volume Outcome Estimates, Stronger Balanced]{This table gives volume outcome coefficient estimates for the treatment effect following a shooting event both overall using equation \ref{eq1_DiD} in Panel A and by scrutiny level using equation \ref{eq3_DiDiD} in Panel B.}
		\label{full_DiD_volumes_A04}
	\end{center}
\end{table}

\begin{table}
	\begin{center}
		\input{C:/Users/david/coding/02_TRAFFIC_BIAS_paper/6000_LaTeX/tables/DiD_rates_A04_full.tex}
		\caption[Rate Outcome Estimates, Stronger Balance]{This table gives rate outcome coefficient estimates for the treatment effect following a shooting event both overall using equation \ref{eq1_DiD} in Panel A and by scrutiny level using equation \ref{eq3_DiDiD} in Panel B.}
		\label{full_DiD_rates_A04}
	\end{center}
\end{table}

%%%%%%%%%%%%%%%%%%%%%%%%%%%%%%%%%%%%%%%%%%%%%%%%%%
%	SINGLE-EVENT DEPARTMENTS 
%%%%%%%%%%%%%%%%%%%%%%%%%%%%%%%%%%%%%%%%%%%%%%%%%%
\clearpage
\subsection{Single-Event Departments}
To evaluate whether my results are driven by multi-event departments, I run my regressions using only single-event departments, which I only observe committing one fatal shooting of an unarmed black subject within the time period, finding consistent results. \\

%	Full DiD + DiDiD Tables
\begin{table}
	\begin{center}
		\input{C:/Users/david/coding/02_TRAFFIC_BIAS_paper/6000_LaTeX/tables/DiD_volumes_A05_full.tex}
		\caption[Volume Outcome Estimates, Single-Event Departments]{This table gives volume outcome coefficient estimates for the treatment effect following a shooting event both overall using equation \ref{eq1_DiD} in Panel A and by scrutiny level using equation \ref{eq3_DiDiD} in Panel B.}
		\label{full_DiD_volumes_A05}
	\end{center}
\end{table}

\begin{table}
	\begin{center}
		\input{C:/Users/david/coding/02_TRAFFIC_BIAS_paper/6000_LaTeX/tables/DiD_rates_A05_full.tex}
		\caption[Rate Outcome Estimates, Single-Event Departments]{This table gives rate outcome coefficient estimates for the treatment effect following a shooting event both overall using equation \ref{eq1_DiD} in Panel A and by scrutiny level using equation \ref{eq3_DiDiD} in Panel B.}
		\label{full_DiD_rates_A05}
	\end{center}
\end{table}

%%%%%%%%%%%%%%%%%%%%%%%%%%%%%%%%%%%%%%%%%%%%%%%%%%
%	MULTI-EVENT DEPARTMENTS 
%%%%%%%%%%%%%%%%%%%%%%%%%%%%%%%%%%%%%%%%%%%%%%%%%%
\clearpage
\subsection{Multi-Event Departments}
Opposite to the previous robustness check, I use the sample of only multi-event departments to see if effects are driven by single-event departments, finding similar estimates. \\

%	Full DiD + DiDiD Tables
\begin{table}
	\begin{center}
		\input{C:/Users/david/coding/02_TRAFFIC_BIAS_paper/6000_LaTeX/tables/DiD_volumes_A06_full.tex}
		\caption[Volume Outcome Estimates, Multi-Event Departments]{This table gives volume outcome coefficient estimates for the treatment effect following a shooting event both overall using equation \ref{eq1_DiD} in Panel A and by scrutiny level using equation \ref{eq3_DiDiD} in Panel B.}
		\label{full_DiD_volumes_A06}
	\end{center}
\end{table}

\begin{table}
	\begin{center}
		\input{C:/Users/david/coding/02_TRAFFIC_BIAS_paper/6000_LaTeX/tables/DiD_rates_A06_full.tex}
		\caption[Rate Outcome Estimates, Multi-Event Departments]{This table gives rate outcome coefficient estimates for the treatment effect following a shooting event both overall using equation \ref{eq1_DiD} in Panel A and by scrutiny level using equation \ref{eq3_DiDiD} in Panel B.}
		\label{full_DiD_rates_A06}
	\end{center}
\end{table}

%%%%%%%%%%%%%%%%%%%%%%%%%%%%%%%%%%%%%%%%%%%%%%%%%%
%	ONLY STATEWIDE DEPARTMENTS 
%%%%%%%%%%%%%%%%%%%%%%%%%%%%%%%%%%%%%%%%%%%%%%%%%%
\clearpage
\subsection{Only Statewide Departments}
With concerns that results may be driven by statewide or municipal departments, I run my regressions on only the sub-sample of statewide departments, such as highway patrols and state troopers, with similar findings. \\

%	Full DiD + DiDiD Tables
\begin{table}
	\begin{center}
		\input{C:/Users/david/coding/02_TRAFFIC_BIAS_paper/6000_LaTeX/tables/DiD_volumes_A07_full.tex}
		\caption[Volume Outcome Estimates, Statewide Departments Only]{This table gives volume outcome coefficient estimates for the treatment effect following a shooting event both overall using equation \ref{eq1_DiD} in Panel A and by scrutiny level using equation \ref{eq3_DiDiD} in Panel B.}
		\label{full_DiD_volumes_A07}
	\end{center}
\end{table}

\begin{table}
	\begin{center}
		\input{C:/Users/david/coding/02_TRAFFIC_BIAS_paper/6000_LaTeX/tables/DiD_rates_A07_full.tex}
		\caption[Rate Outcome Estimates, Statewide Departments Only]{This table gives rate outcome coefficient estimates for the treatment effect following a shooting event both overall using equation \ref{eq1_DiD} in Panel A and by scrutiny level using equation \ref{eq3_DiDiD} in Panel B.}
		\label{full_DiD_rates_A07}
	\end{center}
\end{table}

%%%%%%%%%%%%%%%%%%%%%%%%%%%%%%%%%%%%%%%%%%%%%%%%%%
%	WITHOUT STATEWIDE DEPARTMENTS 
%%%%%%%%%%%%%%%%%%%%%%%%%%%%%%%%%%%%%%%%%%%%%%%%%%
\clearpage
\subsection{Without Statewide Departments}
Finally, I run my regressions excluding statewide departments, once again with generally consistent results. \\

%	Full DiD + DiDiD Tables
\begin{table}
	\begin{center}
		\input{C:/Users/david/coding/02_TRAFFIC_BIAS_paper/6000_LaTeX/tables/DiD_volumes_A08_full.tex}
		\caption[Volume Outcome Estimates, Without Statewide Departments]{This table gives volume outcome coefficient estimates for the treatment effect following a shooting event both overall using equation \ref{eq1_DiD} in Panel A and by scrutiny level using equation \ref{eq3_DiDiD} in Panel B.}
		\label{full_DiD_volumes_A08}
	\end{center}
\end{table}

\begin{table}
	\begin{center}
		\input{C:/Users/david/coding/02_TRAFFIC_BIAS_paper/6000_LaTeX/tables/DiD_rates_A08_full.tex}
		\caption[Rate Outcome Estimates, Without Statewide Departments]{This table gives rate outcome coefficient estimates for the treatment effect following a shooting event both overall using equation \ref{eq1_DiD} in Panel A and by scrutiny level using equation \ref{eq3_DiDiD} in Panel B.}
		\label{full_DiD_rates_A08}
	\end{center}
\end{table}

%%%%%%%%%%%%%%%%%%%%%%%%%%%%%%%%%%%%%%%%%%%%%%%%%%
%	Chapter 2 Appendix
%%%%%%%%%%%%%%%%%%%%%%%%%%%%%%%%%%%%%%%%%%%%%%%%%%
\chapter{Chapter 2 Appendix}
\graphicspath{{C:/Users/david/coding/01_PRISON_PHONE_paper/6000_LaTeX/}}
\addcontentsline{lof}{chapter}{\numberline{\thechapter}Chapter 2 Appendix}
\addcontentsline{lot}{chapter}{\numberline{\thechapter}Chapter 2 Appendix}

\subsection{Robustness Check: Balanced Panel}
I explore the robustness of my null results by considering a balanced panel, finding similar null results. As discussed in the main body of Chapter 2, it is possible that utilizing a larger dataset for my sexual assault outcomes or phone rate changes could yield more precise estimates, or a new outcome altogether with higher frequency could prove useful as well. 

\begin{figure}[H]
	\begin{center}
		\includegraphics[width=4.5in]{inmate_noncon_all_balanced.eps}
		\caption[Alleged Inmate Non-Consensual Sexual Acts Event Study, Balanced]{As in figure \ref{inmate_noncon_all}, there does not appear to be any significant treatment effect for Inmate-on-Inmate Non-Consensual Sexual Act Allegations.}
		\label{inmate_noncon_all_balanced}
	\end{center}
\end{figure}

\begin{figure}[H]
	\begin{center}
		\includegraphics[width=4.5in]{staff_miscond_all_balanced.eps}
		\caption[Alleged Staff Sexual Misconduct Event Study, Balanced]{As in figure \ref{staff_miscond_all}, there does not appear to be any significant treatment effect for Staff-on-Inmate Sexual Misconduct Allegations.}
		\label{staff_miscond_all_balanced}
	\end{center}
\end{figure}

%%%%%%%%%%%%%%%%%%%%%%%%%%%%%%%%%%%%%%%%%%%%%%%%%%
%	Chapter 3 Appendix
%%%%%%%%%%%%%%%%%%%%%%%%%%%%%%%%%%%%%%%%%%%%%%%%%%
\chapter{Chapter 3 Appendix}
\graphicspath{{C:/Users/david/coding/03_GAS_TAX_HETEROGENEITY_paper/6000_LaTeX/}}
\addcontentsline{lof}{chapter}{\numberline{\thechapter}Chapter 3 Appendix}
\addcontentsline{lot}{chapter}{\numberline{\thechapter}Chapter 3 Appendix}

%%%%%%%%%%%%%%%%%%%%%%%%%%%%%%%%%%%%%%%%%%%%%%%%%%
%	BOOTSTRAPPED STANDARD ERRORS
%%%%%%%%%%%%%%%%%%%%%%%%%%%%%%%%%%%%%%%%%%%%%%%%%%
\subsection{Bootstrapped Standard Errors}
As a result of using estimated elasticities as outcomes in my regressions, which introduces uncertainty to our regressions, we correct the standard errors via bootstrapping, running 500 iterations of our main regressions with re-sampling by locale. While this correction does not largely change our results, we do find stronger predictive power from weather quality, border population share and licensure rates, and we find that urbanization and population density may be less explanatory than previous results suggested, although still statistically significant. \\

%	Bootstrap-Corrected Main Predictors, Weighted
\begin{table}[H]
	\begin{center}
		\input{C:/Users/david/coding/03_GAS_TAX_HETEROGENEITY_paper/6000_LaTeX/etable_y1_main_w_boot.tex}
		\label{main_estimates_w_boot}
		\caption[Bootstrap-Corrected Main Predictor Estimates, Weighted]{Using bootstrap-corrected standard errors, we find similar significance for the estimates in columns (1), (2), (3), and (5), while population density becomes less statistically significant, and the weather index becomes the most explanatory in the full regression with all 5 predictors.}
	\end{center}
\end{table}

%	Bootstrap-Corrected Environmental Predictors, Weighted
\begin{table}[H]
	\begin{center}
		\input{C:/Users/david/coding/03_GAS_TAX_HETEROGENEITY_paper/6000_LaTeX/etable_y1_environmental_w_boot.tex}
		\label{environmental_estimates_w_boot}
		\caption[Bootstrap-Corrected Environmental Predictor Estimates, Weighted]{Using bootstrap-corrected standard errors, our results remain unchanged when considering environmental predictors.}
	\end{center}
\end{table}

%	Bootstrap-Corrected Transportation Infrastructure Predictors, Weighted
\begin{table}[H]
	\begin{center}
		\input{C:/Users/david/coding/03_GAS_TAX_HETEROGENEITY_paper/6000_LaTeX/etable_y1_transportation_w_boot.tex}
		\label{transportation_estimates_w_boot}
		\caption[Bootstrap-Corrected Transportation Infrastructure Predictor Estimates, Weighted]{Using bootstrap-corrected standard errors, our results remain largely unchanged when considering transportation network predictors, although the level of urbanization becomes less statistically significant.}
	\end{center}
\end{table}

%	Bootstrap-Corrected Demographic Predictors, Weighted
\begin{table}[H]
	\begin{center}
		\input{C:/Users/david/coding/03_GAS_TAX_HETEROGENEITY_paper/6000_LaTeX/etable_y1_demographic_w_boot.tex}
		\label{demographic_estimates_w_boot}
		\caption[Bootstrap-Corrected Demographic Predictor Estimates, Weighted]{Using bootstrap-corrected standard errors, our results change slightly when considering demographics, with population density becoming less explanatory, and border population share and licensure rate becoming more predictive.}
	\end{center}
\end{table}

%	Bootstrap-Corrected Pop-Share Predictors, Weighted
\begin{table}[H]
	\begin{center}
		\input{C:/Users/david/coding/03_GAS_TAX_HETEROGENEITY_paper/6000_LaTeX/etable_y1_pop_shares_w_boot.tex}
		\label{pop_share_estimates_w_boot}
		\caption[Bootstrap-Corrected Cheaper Gas Border Predictor Estimates, Weighted]{Using bootstrap-corrected standard errors, our results are largely unchanged when considering border population share factors, with border population share becoming more strongly predictive of elasticity.}
	\end{center}
\end{table}

%%%%%%%%%%%%%%%%%%%%%%%%%%%%%%%%%%%%%%%%%%%%%%%%%%
%	UNWEIGHTED ESTIMATES
%%%%%%%%%%%%%%%%%%%%%%%%%%%%%%%%%%%%%%%%%%%%%%%%%%
\subsection{Unweighted Estimates}
As described in the main body of the paper, the number of drivers licenses in each jurisdiction in 1980 are used to analytically weight elasticity regressions, acting as indirect measures of the demand side of each locale's fuel market. To alleviate concerns with this decision, the form of the following tables of estimates are identical to the main results, but without weights applied. We find that the direction of our estimates remains consistent, despite changes in significance, giving credibility to the weighting decision. \\

%	Main Predictors, Unweighted
\begin{table}[H]
	\begin{center}
		\input{C:/Users/david/coding/03_GAS_TAX_HETEROGENEITY_paper/6000_LaTeX/etable_y1_main.tex}
		\label{main_estimates}
		\caption[Main Predictor Estimates, Unweighted]{To check the robustness of our analytic regression weights, we re-run each regression without weights. While the significance of our estimates changes, the signs remain invariant}
	\end{center}
\end{table}

%	Environmental Predictors, Unweighted
\begin{table}[H]
	\begin{center}
		\input{C:/Users/david/coding/03_GAS_TAX_HETEROGENEITY_paper/6000_LaTeX/etable_y1_environmental.tex}
		\label{environmental_estimates}
		\caption[Environmental Predictor Estimates, Unweighted]{Continuing as described in table \ref{main_estimates}, we find largely similar results in our unweighted regressions.}
	\end{center}
\end{table}

%	Transportation Infrastructure Predictors, Unweighted
\begin{table}[H]
	\begin{center}
		\input{C:/Users/david/coding/03_GAS_TAX_HETEROGENEITY_paper/6000_LaTeX/etable_y1_transportation.tex}
		\label{transportation_estimates}
		\caption[Transportation Infrastructure Predictor Estimates, Unweighted]{Continuing as described in table \ref{main_estimates}, we find largely similar results in our unweighted regressions.}
	\end{center}
\end{table}

%	Demographic Predictors, Unweighted
\begin{table}[H]
	\begin{center}
		\input{C:/Users/david/coding/03_GAS_TAX_HETEROGENEITY_paper/6000_LaTeX/etable_y1_demographic.tex}
		\label{demographic_estimates}
		\caption[Demographic Predictor Estimates, Unweighted]{Continuing as described in table \ref{main_estimates}, we find largely similar results in our unweighted regressions.}
	\end{center}
\end{table}

%	Pop-Share Predictors, Unweighted
\begin{table}[H]
	\begin{center}
		\input{C:/Users/david/coding/03_GAS_TAX_HETEROGENEITY_paper/6000_LaTeX/etable_y1_pop_shares.tex}
		\label{pop_share_estimates}
		\caption[Cheaper Gas Border Predictor Estimates, Unweighted]{Continuing as described in table \ref{main_estimates}, we find largely similar results in our unweighted regressions.}
	\end{center}
\end{table}

%%%%%%%%%%%%%%%%%%%%%%%%%%%%%%%%%%%%%%%%%%%%%%%%%%
%	WITH CONTROLS, WEIGHTED ESTIMATES
%%%%%%%%%%%%%%%%%%%%%%%%%%%%%%%%%%%%%%%%%%%%%%%%%%
\subsection{Robustness Check, Elasticity\textsuperscript{2}, Weighted Estimates}
To further explore the robustness of our results, we consider alternate estimates of elasticity, also taken from Bates \& Kim (JAE 2024), taken from their Model 2, which utilizes a series of control variables. We first check our results in weighted regressions below, and in a later section we continue with another check without weights. \\

Our main findings maintain the same sign despite reduced significance, somewhat strengthening our findings and providing evidence against over-sensitivity. \\

%	Main Predictors, Weighted
\begin{table}[H]
	\begin{center}
		\input{C:/Users/david/coding/03_GAS_TAX_HETEROGENEITY_paper/6000_LaTeX/etable_y2_main_w.tex}
		\label{main_estimates_y2_w}
		\caption[Alternate Main Predictor Estimates, Weighted]{Despite using different elasticity estimates, our findings maintain the same direction, and population density remains as the strongest predictor of gas price elasticity.}
	\end{center}
\end{table}

%	Environmental Predictors, Weighted
\begin{table}[H]
	\begin{center}
		\input{C:/Users/david/coding/03_GAS_TAX_HETEROGENEITY_paper/6000_LaTeX/etable_y2_environmental_w.tex}
		\label{environmental_estimates_y2_w}
		\caption[Alternate Environmental Predictor Estimates, Weighted]{Continuing as described in table \ref{main_estimates_y2_w}, we find largely similar results using our alternate elasticity measure.}
	\end{center}
\end{table}

%	Transportation Infrastructure Predictors, Weighted
\begin{table}[H]
	\begin{center}
		\input{C:/Users/david/coding/03_GAS_TAX_HETEROGENEITY_paper/6000_LaTeX/etable_y2_transportation_w.tex}
		\label{transportation_estimates_y2_w}
		\caption[Alternate Transportation Infrastructure Predictor Estimates, Weighted]{Continuing as described in table \ref{main_estimates_y2_w}, we find largely similar results using our alternate elasticity measure.}
	\end{center}
\end{table}

%	Demographic Predictors, Weighted
\begin{table}[H]
	\begin{center}
		\input{C:/Users/david/coding/03_GAS_TAX_HETEROGENEITY_paper/6000_LaTeX/etable_y2_demographic_w.tex}
		\label{demographic_estimates_y2_w}
		\caption[Alternate Demographic Predictor Estimates, Weighted]{Continuing as described in table \ref{main_estimates_y2_w}, we find largely similar results using our alternate elasticity measure.}
	\end{center}
\end{table}

%	Pop-Share Predictors, Weighted
\begin{table}[H]
	\begin{center}
		\input{C:/Users/david/coding/03_GAS_TAX_HETEROGENEITY_paper/6000_LaTeX/etable_y2_pop_shares_w.tex}
		\label{pop_share_estimates_y2_w}
		\caption[Alternate Cheaper Gas Border Predictor Estimates, Weighted]{Continuing as described in table \ref{main_estimates_y2_w}, we find largely similar results using our alternate elasticity measure.}
	\end{center}
\end{table}

%%%%%%%%%%%%%%%%%%%%%%%%%%%%%%%%%%%%%%%%%%%%%%%%%%
%	WITH CONTROLS, UNWEIGHTED ESTIMATES
%%%%%%%%%%%%%%%%%%%%%%%%%%%%%%%%%%%%%%%%%%%%%%%%%%
\subsection{Robustness Check, Elasticity\textsuperscript{2}, Unweighted Estimates}
Continuing in similar fashion to the above, we finally evaluate the sensitivity of our results using both the alternate elasticity measure, $Elasticity^{2}_{s}$ in unweighted regressions. As before, we find largely unchanged signs despite fluctuations in the significance of our estimates. \\

All-in-all, we believe that across these various robustness checks, our relatively-constant results indicate that our analysis is not overly-sensitive to specification changes and serve as a reliable preliminary analysis as to the underlying drivers of heterogeneity in gas price elasticity across U.S. states and D.C. 

%	Main Predictors, Unweighted
\begin{table}[H]
	\begin{center}
		\input{C:/Users/david/coding/03_GAS_TAX_HETEROGENEITY_paper/6000_LaTeX/etable_y2_main.tex}
		\label{main_estimates_y2}
		\caption[Alternate Main Predictor Estimates, Unweighted]{Turning to unweighted estimates using our secondary elasticity outcome, we once again find largely similar results, indicating that our weighting decision is not altering our results unduly.}
	\end{center}
\end{table}

%	Environmental Predictors, Unweighted
\begin{table}[H]
	\begin{center}
		\input{C:/Users/david/coding/03_GAS_TAX_HETEROGENEITY_paper/6000_LaTeX/etable_y2_environmental.tex}
		\label{environmental_estimates_y2}
		\caption[Alternate Environmental Predictor Estimates, Unweighted]{Continuing as described in table \ref{main_estimates_y2}, we find largely similar results in unweighted regressions using our alternate elasticity measure.}
	\end{center}
\end{table}

%	Transportation Infrastructure Predictors, Unweighted
\begin{table}[H]
	\begin{center}
		\input{C:/Users/david/coding/03_GAS_TAX_HETEROGENEITY_paper/6000_LaTeX/etable_y2_transportation.tex}
		\label{transportation_estimates_y2}
		\caption[Alternate Transportation Infrastructure Predictor Estimates, Unweighted]{Continuing as described in table \ref{main_estimates_y2}, we find largely similar results in unweighted regressions using our alternate elasticity measure.}
	\end{center}
\end{table}

%	Demographic Predictors, Unweighted
\begin{table}[H]
	\begin{center}
		\input{C:/Users/david/coding/03_GAS_TAX_HETEROGENEITY_paper/6000_LaTeX/etable_y2_demographic.tex}
		\label{demographic_estimates_y2}
		\caption[Alternate Demographic Predictor Estimates, Unweighted]{Continuing as described in table \ref{main_estimates_y2}, we find largely similar results in unweighted regressions using our alternate elasticity measure.}
	\end{center}
\end{table}

%	Pop-Share Predictors, Unweighted
\begin{table}[H]
	\begin{center}
		\input{C:/Users/david/coding/03_GAS_TAX_HETEROGENEITY_paper/6000_LaTeX/etable_y2_pop_shares.tex}
		\label{pop_share_estimates_y2}
		\caption[Alternate Cheaper Gas Border Predictor Estimates, Unweighted]{Continuing as described in table \ref{main_estimates_y2}, we find largely similar results in unweighted regressions using our alternate elasticity measure.}
	\end{center}
\end{table}
